\section{Analytical Stack}
\label{sec:theory:stack}

A medallion pipeline can be implemented with many software
stacks, ranging from distributed Spark clusters to in-process
single-machine engines. This section introduces the four
classes of tool used in this work --- the analytical database,
the DataFrame runtime, the exploratory notebook environment,
and the visualisation libraries that drive a dashboard surface
--- with a brief description of each class and the specific
implementation chosen.

\subsection{DuckDB and Parquet}
\label{sec:theory:duckdb}

\imageplaceholder{Logo}{Logo of DuckDB}{Download the official logo of DuckDB from https://duckdb.org/; save as images/logos/duckdb.png}

\noindent
DuckDB~\cite{duckdb} is an in-process columnar SQL engine
purpose-built for analytical workloads on a single machine.
Three properties make it well-suited to a medallion
implementation. First, it reads Apache Parquet files
directly: a Parquet file on disk is queryable via
\texttt{read\_parquet('file.parquet')} without an explicit
\texttt{COPY} step. Second, it supports the SQL extensions an
analytical layer needs --- window functions,
\texttt{arg\_max()}, lateral joins, SQL macros for shared
expressions. Third, it has zero deployment overhead: no
separate server, no driver, no network round-trips, which
keeps any solution built on it reproducible from a single
\texttt{pip install}.

\medskip

Parquet is the on-disk format used at the Silver and Gold
layers. The format's columnar layout pairs naturally with the
projection-and-aggregation queries that analytical
applications issue, and ZSTD compression typically reduces raw
CSV by 5--10$\times$ on disk while preserving full query speed.

\subsection{Polars}
\label{sec:theory:polars}

\imageplaceholder{Logo}{Logo of Polars}{Download the official logo of Polars from https://pola.rs/; save as images/logos/polars.png}

\noindent
Polars~\cite{polars} is a high-performance DataFrame library
written in Rust with a Python API. It is the runtime used to
move data out of an analytical database and into application
code. The transfer happens over the Apache Arrow columnar
memory format, so a Polars DataFrame produced from a SQL query
shares memory with the underlying Arrow buffer and incurs no
copy. Polars also exposes a lazy expression API that pushes
filter and projection operations down into the underlying
engine, so the application only materialises the columns it
actually displays.

\subsection{Jupyter Notebook}
\label{sec:theory:jupyter}

\imageplaceholder{Logo}{Logo of Jupyter}{Download the official logo of Jupyter from https://jupyter.org/; save as images/logos/jupyter.png}

\noindent
Jupyter Notebook~\cite{jupyter} is a literate-computing
environment in which code, output, and prose live in the same
document. It is the canonical tool for exploratory data
analysis: the analyst writes a SQL or DataFrame expression,
runs it, inspects the result inline, and refines the next
query against what they just saw. Hypotheses verified
interactively in a notebook can later be promoted into a
materialised view in the Gold layer without rewriting the
underlying SQL. The notebook environment integrates natively
with both DuckDB and Polars.

\subsection{Streamlit}
\label{sec:theory:streamlit}

\imageplaceholder{Logo}{Logo of Streamlit}{Download the official logo of Streamlit from https://streamlit.io/; save as images/logos/streamlit.png}

\noindent
Streamlit~\cite{streamlit} is a Python framework for building
data applications with a reactive programming model. The
application is written as an ordinary Python script;
Streamlit re-executes the script on every user interaction and
renders the resulting widgets and figures to a browser. State
is held in a shared session-state object that survives across
re-executions. The reactive model is well-suited to MCV-style
dashboards (Section~\ref{sec:theory:mcv}) because filter
propagation reduces to writing a value to the session-state
object: every view that reads from that key re-renders
automatically.

\subsection{pydeck}
\label{sec:theory:pydeck}

\imageplaceholder{Logo}{Logo of pydeck}{Download the official logo of pydeck/deck.gl from https://deck.gl/; save as images/logos/pydeck.png}

\noindent
pydeck~\cite{pydeck} is the Python wrapper for deck.gl, a
WebGL-based geospatial visualisation library originally
released by Uber. It draws thousands of map features --- stop
dots, route polylines, catchment polygons --- at interactive
frame rates and exposes click events back to the host
application. The \texttt{ScatterplotLayer} used in this work
supports continuous colour gradients and per-point tooltips,
which are the two primitives a stop-density map requires.

\subsection{Plotly}
\label{sec:theory:plotly}

\imageplaceholder{Logo}{Logo of Plotly}{Download the official logo of Plotly from https://plotly.com/; save as images/logos/plotly.png}

\noindent
Plotly~\cite{plotly} is a charting library that provides the
canonical chart vocabulary --- line, area, bar, scatter,
treemap, histogram, box plot, heatmap --- used by analytical
dashboards. It supports server-side interactivity: click,
brush, and hover events are reported back to the host
application so that a chart click can trigger a state change
elsewhere. Plotly's chart catalogue covers every analytical
question identified in Chapter~\ref{chapter:proposed} without
extension.

\bigskip

The six tools above cover the entire analytical stack from raw
data to rendered dashboard. None is exotic; together they form
a small, maintainable surface area that can be deployed from a
single Python environment with no external services.

Of the thirty-one canonical charts specified in the dashboard
(Section~\ref{sec:proposed:dashboard}), twenty-two render
natively through \texttt{st.plotly\_chart}: heatmaps, ranked
bars, parallel coordinates, treemaps, scatter quadrants, stacked
areas, promotion-lift bars, and the demand-forecast multi-line
chart. Six charts use \texttt{st.pydeck\_chart} for deck.gl
integration: 3.1 network overview, 3.2 boarding density
HexagonLayer, 3.3 stop demand explorer, 3.4 ghost-stop map, 3.5
bus--metro feeder map, and 1.5 school proximity map. The
remaining three charts (6.1 the UMAP route constellation, 6.2
the GroupSet/UpSet archetype matrix, 6.5 the Marey
time--distance diagram), plus the hub landing page, the
scrollytelling Stories layer, and the universal entity-detail
panel (invariant~A3 of Section~\ref{sec:proposed:invariants}),
are injected via \texttt{st.components.v1.html} as self-contained
HTML components that own their own interactivity.
